\documentclass[10pt]{article}

%%% These are some packages that are useful
\usepackage{amsmath,amssymb, amscd,amsbsy, amsthm, enumerate}
\usepackage[export]{adjustbox}
\usepackage{lastpage}
\usepackage[top=1in, bottom=1in, left=1in, right=1in]{geometry}
\usepackage[unicode]{hyperref}
\usepackage{tikz, pgfplots, xcolor, fancyhdr}
\usepackage{multicol,caption}
\usepackage{lipsum}
\usepackage[version=4]{mhchem}
\usepackage{float}
\usepackage{enumitem}
\usepackage{xurl}

%%% Page formatting
%\setlength{\headsep}{30pt}
\setlength{\textheight}{9in}
\newcommand{\tab}{\hspace{1cm}}
%\setlength{\parindent}{25pt}

\title{Know Thy Enemy (Part 5) 	: Astarte (Ishtar, Aphrodite, Hathor, ...)}
\author{Antonius Torode}

%%% Header and Footer Info
\pagestyle{fancy}
\fancyhead[L]{{\large Template - \textbf{Change 03}}}
\fancyhead[C]{\today}
\fancyhead[R]{Name: Antonius Torode}


\fancyhf{} % sets both header and footer to nothing
\renewcommand{\headrulewidth}{0pt}
% your new footer definitions here

\fancyfoot[L]{}
\fancyfoot[C]{}
\fancyfoot[R]{\thepage\ of \pageref{LastPage}}

% Used to define spacing and format of References
\let\OLDthebibliography\thebibliography
\renewcommand\thebibliography[1]{
	\OLDthebibliography{#1}
	\setlength{\parskip}{0pt}
	\setlength{\itemsep}{0pt plus 0.3ex}
}

\newenvironment{Figure}
  {\par\medskip\noindent\minipage{\linewidth}}
  {\endminipage\par\medskip}


%%% Document Starts now
\begin{document}

\maketitle
\thispagestyle{fancy}


%\begin{abstract}
%TODO
%\end{abstract}

\begin{multicols}{2}

In a previous part of this series, I demonstrated that the old gods are still around and active in today's age - as they have been throughout history \cite{part3}. My last message was about Baal\cite{baal}. Throughout history, once Baal takes a foothold among a people, it opens up a floodgate of other entities to step in and take hold. One of those beings, is commonly known as Ishtar. I believe this being to be one of the most obvious growing influences in many parts of the world today (specifically seen over the past few years). So, who is Ishtar, and how are they influencing the world today?

Throughout ancient times, Israel turned their ways to false worship of old gods:
\begin{quotation}
``They forsook the Lord and served Baal and \textit{the Ashtoreths}.'' - Judges 2:13
\end{quotation}
One of the common ones they continually turned to is called \textit{the Ashtoreths}. 
\begin{quotation}
``Then the children of Israel again did evil in the sight of the Lord, and served the Baals and \textit{the Ashtoreths}, the gods of Syria, the gods of Sidon, the gods of Moab, the gods of the people of Ammon, and the gods of the Philistines; and they forsook the Lord and did not serve Him.'' - Judges 10:6
\end{quotation}
The word \textit{Ashtoreths} is ``the collective name of the Canaanite fertility-war goddess Astarte and her images \cite{Ashtoreths strongs}.'' Astarte was also the goddess of the Sidonians\footnote{1 Kings 11:5 - ``For Solomon went after Ashtoreth the goddess of the Sidonians, and after Milcom the abomination of the Ammonites.''}. Astarte is commonly known as the goddess of war and sexual love \cite{Britannica}. This is a being that appeared throughout history in various cultures; known by different names but representing the same ideologies. 

\begin{quotation}
``Her Akkadian counterpart was Ishtar. Later she became assimilated with the Egyptian deities Isis and Hathor (a goddess of the sky and of women), and in the Greco-Roman world with Aphrodite, Artemis, and Juno.'' - Britannica Encyclopedia\cite{Britannica} 
\end{quotation}

Just like how Baal had many names all representing the same entity\cite{baal}, so too is the case with Astarte/Ishtar/Aphrodite/Inanna/etc \cite{ishtar is astarte, ishtar is inanna}. A big difference between the two is that Baal's presence in the modern age is subtle; specifically from the perspective of there not being an obvious Baal-worshiping people. Once you see the patterns and match that with historical trends, you should see it; though it's not in-your-face. With Ishtar, that is not the case. Once Baal has opened the door and taken over a nation, other demonic forces are free to step in and be as obvious and open as they want - they become free to flaunt their true selves in the open. 

Many have noticed a trend over the past decade, where the LGBTQIA2S+\cite{lgbtqia2s} community has been taking a foothold in our culture.  Kids are being convinced at record rates that they are in the wrong gendered body, that they don't fit among the two genders, or other non-biblical nonsense. There has been a seemingly sudden explosion of this type of individual identity which has quadrupled over the past decade or so - raising from 8\% of young women in 2012 to 31\% in 2024\cite{lgbtqrise}. What has changed over the past decade to cause such a great expansion in this ideology? 

There are many reasons for this ideology taking a hold, but an obvious one (most relevant to this message) is the \textit{open} worship of Ishtar. This year, the 11th annual Ishtarfest was celebrated. Although a relatively tiny festival, their purpose is very clear:
\begin{quotation}
``Step into the radiance of the Fertile Crescent and rediscover the ancient wisdom that shaped civilization itself. Ishtarfest is a vibrant, community-centered festival \textit{dedicated} to the Goddess Inanna/Ishtar—Queen of Heaven and Earth, Lady of Love and War, and one of the earliest champions of justice, sovereignty, and sacred order. \cite{Ishtarfest 2026}''
\end{quotation}
There can be no doubt that active Ishtar worship occurs in our modern world today. That fact also ignores anything that potentially occurs behind closed doors. Furthermore, the open and active \textit{worship} of Ishtar corresponds perfectly with the exponential rise in the practices she teaches. 

Similar to how Baal appears throughout scripture, Astarte also appears in many places \textit{without} an explanation or introduction. The reader is assumed to know who or what Astarte is. So, what exactly does Ishtar teach and represent? Ultimately, Astarte is a demonic entity whose mind we cannot truly know\cite{part3}. However, we can learn from historical patterns, myths, and obscure teachings to get a very clear picture of what this entity represents. This allows us to both discern and understand where certain practices come from, and why they should be avoided. Most people who hear of the various things that Ishtar represents will immediately see the connections to what is happening in our cultures today - those connections are often self evident. 

At a high level, Astarte is often characterized by sexual immorality and prostitution \cite{Astarte bible hub}. This sexual immorality manifests itself in many different ways. 

One way this manifests itself today should seem familiar with what's occurring in the culture around us. Inanna is largely known for ``binary breaking and gender bending \cite{Astarte Gender Fluidity}.'' Throughout history, one repeating common sign of her followers is the cultist practice of gender fluidity and transformation. There's an old piece of devotional work by the poet and high priestess Enheduanna (dated circa 2300 BCE) called \textit{The Hymn to Inanna} \cite{Hymn To Innana, Lady of Largest Heart}. A paper published by the University of Pennsylvania Museum looks at this from a historical context and states:
\begin{quotation}
``It may simply have been enough to understand that Inanna is the one ‘turning a male into a female and a female into a male,’ as explained in the Lady of the Largest Heart hymn\footnote{The phrase \textit{Lady of the Largest Heart} is another title and reference to Innana \cite{Lady of Largest Heart}.} \cite{Hymn To Innana}''
\end{quotation}
But they go further to state that Innana's purview isn't just the physical transformation between the genders - but also psychological, social, and more. There's a very clear reason as to why Inanna represents these things.

Ishtar represents the destruction of established boundaries - a being that is against all things that God put in order \cite{opposites, distortion}. In ancient Mesopotamia, she disrupted the social order and distorted the normal boundaries of society \cite{distortion}. In our culture today, she is doing the same thing. 

In the beginning, God created male and female\footnote{Genesis 1:27, Genesis 5:2, Genesis 6:19, Mathew 19:4, Mark 10:6}. That is a clear physical boundary established by God at creation. 




\begin{thebibliography}{9}
{\footnotesize
	
	\bibitem{part3} Torode, Antonius, \textit{Know Thy Enemy (part 3): The Old Gods}, June 2026.  \url{https://github.com/torodean/bible-notes/blob/main/sermonettes/Know%20Thy%20Enemy/The%20Old%20Gods/The%20Old%20Gods.pdf}
		
	\bibitem{baal} Torode, Antonius, \textit{Know Thy Enemy (part 4): Baal}, August 2026.  \url{https://github.com/torodean/bible-notes/blob/main/sermonettes/Know%20Thy%20Enemy/Baal/Baal.pdf}
	
	\bibitem{Ashtoreths strongs} Bible Hub: Ashtoreths. Accessed August 8 2026. \url{https://biblehub.com/hebrew/6252.htm}
	
	\bibitem{Britannica} Britannica Editors, \textit{Astarte}, Encyclopedia Britannica, January 23, 2025, Accessed August 8, 2026. \url{https://www.britannica.com/topic/Astarte-ancient-deity}
	
	\bibitem{ishtar is astarte} Sugimoto, David T., \textit{Transformation of a Goddess: Ishtar - Astarte - Aphrodite}, University of Zurich, 2014.  \url{https://www.scribd.com/document/384308293/Transformation-of-a-Goddess}
	
	\bibitem{ishtar is inanna} Yaǧmur Heffron, \textit{Inana/Ištar (goddess)}, Ancient Mesopotamian Gods and Goddesses, Oracc and the UK Higher Education Academy, 2024, Accessed August 22, 2026. \url{https://oracc.museum.upenn.edu/amgg/Listofdeities/InanaIshtar/index.html} 
	
	\bibitem{Ishtarfest 2026} Hands of Change, \textit{Where Inanna Walks: Memory, Justice, and Sacred Community}, 2026. Accessed August 15, 2026. \url{https://handsofchange.gumroad.com/l/ycapzw}
	
	\bibitem{lgbtqia2s} Tiersa, \textit{LGBTQIA2S+: What Does It All Mean?}, Medicine Hat Public Library, June 29, 2021, Accessed August 15, 2026. \url{https://mhpl.shortgrass.ca/blog/lgbtqia2s-what-does-it-all-mean}
	
	\bibitem{lgbtqrise}  Daniel A. Cox, et. al., \textit{What’s Behind the Rapid Rise in LGBTQ Identity?}, Survey Center On American Life, March 6, 2025, Accessed August 15, 2026. \url{https://www.americansurveycenter.org/newsletter/whats-behind-the-rapid-rise-in-lgbtq-identity/}
		
	\bibitem{Astarte bible hub} Bible Hub: Astarte. Accessed August 22 2026. \url{https://biblehub.com/topical/a/astarte.htm}
	
	\bibitem{Astarte Gender Fluidity}  Pareja, M.N. \textit{Potnia’s Participants: Considering the Gala, Assinnu, and Kurgarrû in an Aegean Context}, Arts 2024, 13, 20, January 24 2024. \url{https://www.mdpi.com/2076-0752/13/1/20}
	
	\bibitem{Hymn To Innana} Mark, Joshua J., \textit{Hymn to Inanna: An Ancient Praise Poem}, World History Encyclopedia, March 17, 2026, Accessed August 29, 2026.  \url{https://www.worldhistory.org/article/2109/hymn-to-inanna/}
	
	\bibitem{Lady of Largest Heart} Betty De Shong Meador, \textit{Inanna, Lady of Largest Heart: Poems of the Sumerian High Priestess Enheduanna}, Book, February 2001. \url{https://utpress.utexas.edu/9780292752429/}
	
	\bibitem{opposites} Rivkah Harris, \textit{Innana-Ishtar As Paradox And A Coincidence Of Opposites}, The University Of Chicago, History of Religions, Volume 30, Number 3, February 1, 1991. \url{https://www.journals.uchicago.edu/doi/abs/10.1086/463228}
	
	\bibitem{distortion} Leana Wessels, \textit{An analysis of the extent to which the trickster archetype can be applied to the goddess Inanna / Ishtar}, Sabinet African Journals, January 1, 2013. \url{https://journals.co.za/doi/10.10520/EJC139821}
}
\end{thebibliography}

\end{multicols}


%%%%%%%%%%%%%%%%%%%%%%%%%%%%%%%%%%%%%%%%%%%%%%%%%%%%%%%%%%%%%%%%%%%%%%%%%%%%%%%%%%%%%%%%%%%
\end{document}





















}{den}